% Source: https://artofproblemsolving.com/community/c3h1632664_mathcounts_question_not_problem
% Stable identity: aops-topic-1632664-q01
\textbf{(a)} The roots are $1$ (once) and $2$ (twice, since $(n-2)^2$ is a factor). Counted with multiplicity the roots are $1,2,2$, with sum $\boxed{5}$; the distinct roots are $1,2$, with sum $\boxed{3}$.

\textbf{(b)} By Vieta's formulas, for a monic cubic $n^3+pn^2+qn+r=(n-r_1)(n-r_2)(n-r_3)$ the sum $r_1+r_2+r_3=-p$, where $r_1,r_2,r_3$ is the list of roots \emph{with multiplicity}: here $-(-5)=5$. This is why ``the sum of the roots'' normally means the sum with multiplicity, i.e.\ $5$. A problem that wants each root once usually says ``distinct roots'', and a careful problem states which is meant.
